problem:  
Let \((X, d, \mathfrak{m}, p)\) be a non-collapsed \(RCD(K,N)\) space arising as the measured Gromov–Hausdorff limit of a sequence of \(n\)-dimensional Riemannian manifolds \((M_i,g_i,p_i)\) with \(\mathrm{Ric}_{M_i} \ge K\) and \(\operatorname{vol}(B_1(p_i)) \ge v_0 > 0\).  
For a point \(x \in X\), denote by \(\mathcal{T}_x X\) the (possibly nonunique) set of tangent cones at \(x\), and let  
\[
k(x) = \max \left\{ \, k \in \mathbb{N} \mid \exists\, T \in \mathcal{T}_x X \text{ with a factor } \mathbb{R}^k \text{ in its metric splitting } \right\}.
\]  
It is known that \(\dim_{\mathrm{Haus}}(X) = \sup_x k(x)\).  

Open questions:  
1. Is the function \(x \mapsto k(x)\) upper semicontinuous on \(X\) in the general \(RCD(K,N)\) setting, even when tangent cones are not unique?  
2. If there exists an open subset \(U \subset X\) such that all tangent cones at points in \(U\) admit an isometric \(\mathbb{R}^k\)-factor for the same maximal \(k\), does \(U\) necessarily carry a \(C^{1,\alpha}\)-Riemannian manifold structure whose metric tensor coincides (locally in measure) with that of the ambient \(RCD\) metric?  

Why is it a "good" problem:  
This problem probes the yet-unresolved frontier between nonsmooth metric measure geometry and classical Riemannian geometry. While the Cheeger–Colding theory and its extension to \(RCD(K,N)\) spaces established deep structure and rectifiability results, the precise continuity behavior of the maximal Euclidean factor dimension function \(k(x)\) across the space remains unknown, especially without uniqueness of tangent cones. The second part asks whether geometric “flatness in all tangent cones” implies an actual manifold structure—seeking a bridge between infinitesimal geometry and macroscopic differentiability in the \(RCD\) realm.  

Conceptually, this is a "narrow entrance, profound depth" question: the statement is simple, yet answering it would require major advances in understanding stability of splitting, quantitative cone rigidity, and fine regularity under synthetic Ricci curvature bounds. It unites metric geometry, analysis on metric spaces, and geometric measure theory, extending the Cheeger–Colding paradigm into the synthetic setting, making it a genuinely open, research-level, and profoundly integrative problem.